% Generated draft: requires mathematical and source review.
\documentclass[11pt]{article}
\usepackage[margin=25mm]{geometry}
\usepackage{fontspec}
\setmainfont{Latin Modern Roman}
\usepackage{amsmath,amssymb,mathtools,unicode-math}
\setmathfont{Latin Modern Math}
\usepackage{enumitem,needspace,xurl,hyperref}
\setlength{\parindent}{0pt}
\setlength{\parskip}{5pt}
\setlength{\emergencystretch}{3em}
\newcommand{\N}{\mathbb{N}}
\newcommand{\Z}{\mathbb{Z}}
\newcommand{\Q}{\mathbb{Q}}
\newcommand{\R}{\mathbb{R}}
\newcommand{\C}{\mathbb{C}}
\newcommand{\F}{\mathbb{F}}
\newcommand{\A}{\mathbb{A}}
\newcommand{\PP}{\mathbb{P}}
\newcommand{\E}{\mathbb{E}}
\newcommand{\eps}{\varepsilon}
\newcommand{\dd}{\,\mathrm{d}}
\begin{document}
\section*{UnsolvedMath ID: 20003095}
\textbf{Source problem code:} catalogue entry
\subsection{Definitions and mathematical statement}
In finite-type Teichm\"uller space, let $G(\lambda)=(X_t)$ be a Thurston stretch line directed by a maximal geodesic lamination $\lambda$. K. Rafi asks how this geodesic behaves and whether a large projection to the arc-and-curve graph of a subsurface $Y$ forces an interval on which every component of $\partial Y$ has bounded hyperbolic length, together with a suitable graph-geodesic comparison. The source literally makes the final comparison about a boundary component $\beta$. Such a component is peripheral in $Y$ and is not an essential vertex of its usual arc-and-curve graph. The aid explicitly replaces the subject by the projected short-marking shadow; that is a substantive interpretation, not an OCR repair.

For a hyperbolic structure $X_t$, choose a short marking $\mu_t$ and let $\pi_Y(\mu_t)$ be its subsurface projection to $\mathcal{AC}(Y)$. The shadow is the resulting coarse path. Write $d_Y(X_s,X_t)$ for its endpoint projection coefficient, with a fixed convention for bounded projection sets and choice errors. The original question does not specify that convention, thresholds for large or bounded, the subsurface range, or constant dependencies.

The supplied aid invokes a Lenzhen--Modami--Rafi--Tao active-interval theorem for a finite Thurston-geodesic segment $(X_t)_{t\in[a,b]}$ and a non-annular subsurface $Y$ intersecting its maximally stretched lamination. In the applicable theorem range it reports an active interval
\[
 J_Y=[c,d]\subseteq[a,b]
\]
on which each $\beta\subset\partial Y$ has uniformly bounded length, and the short-marking shadow in $\mathcal{AC}(Y)$ is a reparametrized quasi-geodesic. This is a coarse graph statement, not a claim that $\beta$ itself is close to an essential graph geodesic or that the original time parameter has a fixed quasi-geodesic speed. The precise intersection condition, subsurface complexity and uniformity hypotheses require the omitted reference.

If projected movement on each outside segment $[a,c]$ and $[d,b]$ is bounded by $D$, the aid's quantitative localization in its stipulated coefficient convention is
\[
 \bigl|d_Y(X_a,X_b)-d_Y(X_c,X_d)\bigr|\leq2D+4.
\]
It also reports that endpoint projection greater than $2D+4$ forces a nontrivial active interval. The additive $4$ records bounded-choice/projection errors in the report's convention and is not asserted invariant under every alternate definition of $d_Y$. In particular, degenerate-interval and same-time coefficients must be normalized consistently before applying that threshold. Annular subsurfaces are not included by this stated theorem.

For the first part of the question, the summary describes a stretch line as exponential in the shearing coordinates for its directing completion, schematically
\[
 \sigma_\lambda(X_t)=e^{t-t_0}\sigma_\lambda(X_{t_0}),
\]
and identifies its maximally stretched lamination with the largest chain-recurrent part of that completion. The exact coordinate domain and lamination conventions are not supplied. These geometric descriptions and the active-interval statement are reported partial answers, not a classification of every possible stretch-line shadow. The trailing dots in the extraction are page layout; the defective grammar and the final $\beta$ are explicitly reported to occur in the PDF itself.
\subsection{Short English statement}
The proposed interval concerns short boundary lengths and the projected marking path, since a boundary component is not an essential vertex of its own subsurface graph. The aid reports a quantitative active-interval theorem for non-annular subsurfaces, with constants tied to its projection convention, and an exponential shearing description of stretch lines.
\subsection{Sources}
\begin{itemize}
\item Local UnsolvedMath registry, permanent ID 20003095, source code catalogue entry; \texttt{aim-topology-notes.json}, notes[182].
\item AIM, \emph{Lipschitz metric on Teichmueller space}, Problem 2.5, K. Rafi: \url{https://aimath.org/pastworkshops/lipschitzteichproblems.pdf}. The supplied verification says the final $\beta$ and grammatical defect are in the PDF; trailing dots are layout. This draft does not independently inspect it.
\item Lenzhen--Modami--Rafi--Tao active-interval result and stretch/shearing-coordinate facts, named in the supplied research summary without complete bibliography or theorem hypotheses.
\item Machine-generated research attempt 1: partial result, partially solved, rigor labelled ``proved,'' low novelty confidence. Not peer reviewed; independently verify before citation. Registry updated 14 August 2026.
\end{itemize}
\end{document}
